\begin{onehalfspace}

Questa sezione riporta la campagna parametrica schema~2 sulla
strategia TOP-K, di cui il Capitolo~\ref{chap:risultati} riporta la sintesi. Il protocollo usa $N=15$, $a=7$,
$n_{\mathrm{count}}=8$, 30 repliche, budget massimo 50 e, salvo diversa
indicazione, $1\,024$ shot. Il preset di riferimento è
$\lambda_{1q}=10^{-3}$, $\lambda_{2q}=10^{-2}$,
$T_1/T_2=100/80\,\mu$s e readout simmetrico $2\%$.

M1 e TOP-K condividono lo stesso istogramma per ogni coppia
replica--iterazione. I fallimenti sono codificati con sentinella 51 e i
pareggi con una priorità SHA-256 dipendente dal seed. I confronti usano
Wilcoxon appaiato--Pratt. Poiché non è applicata una correzione globale
fra le otto famiglie, i $p$ di questa sezione hanno finalità
esplorativa.

Il materiale seguente approfondisce le Sezioni~\ref{sec:setup_m1}--\ref{sec:parametri_mitigazione}
del corpo. In quei punti restano il disegno, i risultati decisivi e la loro interpretazione;
qui sono raccolti il contratto riproducibile, le validazioni accessorie e tutte le diagnostiche
necessarie a verificarli senza interrompere il filo della discussione.

% -----------------------------------------------
\subsection{Contratto sperimentale e perimetro delle validazioni}
\label{app:fase1_contratto}

Il percorso rumoroso è limitato a $N=15$, $a=7$ e otto qubit di
controllo. Smolin, Smith e Vargo hanno mostrato che una versione compilata
di Shor costruita sfruttando la conoscenza della risposta può fattorizzare
prodotti di due primi in tempo costante e con due soli qubit coerenti;
l'istanza non è quindi presentata come evidenza di scalabilità, ma come
banco di prova comune per strategie di post-processing sul medesimo
istogramma~\cite{smolin2013}.

Il circuito è compilato nella base \texttt{rz/sx/x/cx} con
\texttt{optimization\_level=2} e seed del transpiler 20260819. I gate RZ
sono virtuali; SX/X hanno durata illustrativa di 50\,ns e CX di 300\,ns.
Il modello uniforme combina depolarizzazione a uno e due qubit,
rilassamento/dephasing e readout: non riproduce una calibrazione né la
topologia di una QPU specifica. Ogni iterazione usa $1\,024$ shot, il budget
massimo è 50 e ogni confronto principale comprende 30 repliche. Le
strategie della stessa replica ricevono il medesimo istogramma; i pareggi
sono risolti mediante una priorità SHA-256 dipendente dal seed. I fallimenti
al budget sono codificati con sentinella 51 e le differenze appaiate sono
valutate con Wilcoxon--Pratt e correzione di Holm entro ciascuna famiglia.

Una campagna storica usava semi consecutivi che condividevano quasi tutti
gli shot derivati da Aer. I relativi valori non sono riutilizzati. La
campagna schema~2 separa replica e iterazione di più del budget di shot e
registra nel JSON seed, manifest, revisione del rumore e revisione del
post-processing.

\begin{table}[H]
\centering
\caption[Contratto riproducibile della campagna N=15]{Invarianti del
circuito compilato impiegato da training, baseline, sweep e ZNE.}
\label{tab:correzione_seed}
\label{tab:circuit_depth}
\begin{tabular}{lr}
\toprule
\textbf{Grandezza} & \textbf{Valore} \\
\midrule
Qubit totali & 12 \\
Profondità / size & $412/604$ \\
RZ / SX / CX & $302/70/224$ \\
Misure & 8 \\
SHA-256 circuito & \texttt{c7f4cf3bc1735dce9d5b56d8e09686e...} \\
Qiskit / Aer & $2.5.0/0.17.2$ \\
\bottomrule
\end{tabular}
\end{table}

L'aritmetica Beauregard per $N=21$ e $N=35$ è stata verificata
separatamente nel caso ideale mediante truth table esaustive dei
moltiplicatori controllati, pulizia delle ancille e confronto della legge
QPE con quella teorica. La distanza di variazione totale è
$1.38\times10^{-9}$ per $N=21$, $a=2$, $r=6$ e
$2.50\times10^{-13}$ per $N=35$, $a=6$, $r=2$. Per $N=21$, otto
controlli e compilazione congiunta nella base esecutiva producono 21\,036
CX e profondità 23\,081. Al valore diagnostico
$\lambda_{2q}=10^{-3}$, la probabilità che nessuna CX applichi un Pauli
non-identità è
\[
  \left(1-\frac{15}{16}\lambda_{2q}\right)^{21036}
  =2.70\times10^{-9};
\]
non è una fedeltà né una probabilità di fattorizzazione.

\begin{table}[H]
\centering
\caption[Perimetro delle validazioni]{Separazione fra validazione ideale
e campagna rumorosa. Nel corpo è mantenuto il solo confronto centrale su
$N=15$ e il controllo di scala su $N=21$; $N=35$ resta documentato qui.}
\label{tab:costo_simulazione}
\begin{tabular}{lcc}
\toprule
\textbf{Istanza} & \textbf{Ideale} & \textbf{Rumorosa} \\
\midrule
$N=15$, $a=7$ & truth table e QPE & campagna completa \\
$N=21$, $a=2$ & truth table e QPE & esclusa per costo \\
$N=35$, $a=6$ & truth table e QPE & esclusa per costo \\
\bottomrule
\end{tabular}
\end{table}

% -----------------------------------------------
\subsection{Variazione di K}
\label{subsec:sweep_k}

\begin{table}[H]
\centering
\caption[Sweep del numero di candidati]{M1 e TOP-K sui medesimi 30
vettori di iterazioni. $K=1$ coincide con il controllo M1.}
\label{tab:sweep_k}
\begin{tabular}{cccccc}
\toprule
$K$ & $\bar M_1$ & $\bar M_{\mathrm{TOP}K}$ & successo & $\rho$ & $p$ \\
\midrule
1 & 1.37 & 1.37 & $100\%$ & 1.000 & 1.000 \\
2 & 1.37 & 1.00 & $100\%$ & 1.367 & $<0.001$ \\
3 & 1.37 & 1.00 & $100\%$ & 1.367 & $<0.001$ \\
4 & 1.37 & 1.00 & $100\%$ & 1.367 & $<0.001$ \\
6 & 1.37 & 1.00 & $100\%$ & 1.367 & $<0.001$ \\
8 & 1.37 & 1.00 & $100\%$ & 1.367 & $<0.001$ \\
\bottomrule
\end{tabular}
\end{table}

La soglia osservata è $K=2$, non $K=r=4$. Questo risultato dipende
dalla banda di esiti che la post-elaborazione mappa su un periodo utile e
non costituisce una regola generale per altre istanze.

% -----------------------------------------------
\subsection{Variazione di
\texorpdfstring{$\lambda_{2q}$}{lambda 2q}}
\label{subsec:sweep_eps}

Per il canale Aer a due qubit,
$\mathcal E(\rho)=(1-\lambda)\rho+\lambda I/4$, la probabilità
aggregata della componente Pauli non-identità è $15\lambda/16$. Sul
circuito da 224 CX il proxy di nessun evento è quindi
$(1-15\lambda_{2q}/16)^{224}$; non è una probabilità di successo.

\begin{table}[H]
\centering
\caption[Sweep del parametro Aer $\lambda_{2q}$]{Il tasso di successo è
$100\%$ per tutte le 30 repliche di ciascun punto.}
\label{tab:sweep_eps}
\label{tab:sweep_eps2q}
\begin{tabular}{cccccc}
\toprule
$\lambda_{2q}$ & proxy & $\bar M_1$ & $\bar M_{\mathrm{TOP4}}$ & $\rho$ & $p$ \\
\midrule
$10^{-3}$ & $8.11\times10^{-1}$ & 1.93 & 1.00 & 1.933 & $<0.001$ \\
$5\times10^{-3}$ & $3.49\times10^{-1}$ & 1.90 & 1.00 & 1.900 & $<0.001$ \\
$10^{-2}$ & $1.21\times10^{-1}$ & 1.37 & 1.00 & 1.367 & $<0.001$ \\
$2\times10^{-2}$ & $1.44\times10^{-2}$ & 1.33 & 1.00 & 1.333 & 0.001 \\
$5\times10^{-2}$ & $2.14\times10^{-5}$ & 1.33 & 1.00 & 1.333 & $<0.001$ \\
$10^{-1}$ & $2.65\times10^{-10}$ & 1.57 & 1.07 & 1.469 & 0.001 \\
\bottomrule
\end{tabular}
\end{table}

Al punto estremo TOP-4 richiede in media 1.07 iterazioni: la robustezza
non è assoluta, pur restando il successo $30/30$ entro il budget.

% -----------------------------------------------
\subsection{Variazione degli Shot}
\label{subsec:sweep_shots}

\begin{table}[H]
\centering
\caption[Sweep degli shot]{Risultati su 30 repliche per punto.}
\label{tab:sweep_shots}
\begin{tabular}{ccccc}
\toprule
shot & $\bar M_1$ & $\bar M_{\mathrm{TOP4}}$ & $\rho$ & $p$ \\
\midrule
128  & 1.37 & 1.00 & 1.367 & 0.001 \\
256  & 1.57 & 1.00 & 1.567 & $<0.001$ \\
512  & 1.57 & 1.00 & 1.567 & 0.001 \\
1\,024 & 1.37 & 1.00 & 1.367 & $<0.001$ \\
2\,048 & 1.70 & 1.00 & 1.700 & $<0.001$ \\
\bottomrule
\end{tabular}
\end{table}

TOP-4 è stabile già a 128 shot; M1 oscilla senza andamento monotono,
coerentemente con una competizione campionaria fra quattro picchi.

% -----------------------------------------------
\subsection{Analisi Congiunta
\texorpdfstring{$K\times\lambda_{2q}$}{K x lambda 2q}}
\label{subsec:sweep_joint}

\begin{table}[H]
\centering
\caption[Griglia $K\times\lambda_{2q}$]{Ogni cella è
$\bar M_{\mathrm{TOP}K}$ su 30 repliche e ha successo $100\%$.}
\label{tab:sweep_joint}
\begin{tabular}{ccccc}
\toprule
$K$ & $5\times10^{-3}$ & $10^{-2}$ & $2\times10^{-2}$ & $5\times10^{-2}$ \\
\midrule
2 & 1.00 & 1.00 & 1.00 & 1.00 \\
4 & 1.00 & 1.00 & 1.00 & 1.00 \\
6 & 1.00 & 1.00 & 1.00 & 1.00 \\
8 & 1.00 & 1.00 & 1.00 & 1.00 \\
\bottomrule
\end{tabular}
\end{table}

% -----------------------------------------------
\subsection{Parametri Secondari}
\label{subsec:sweep_secondari}

RZ è virtuale (0\,ns), SX/X dura 50\,ns e CX 300\,ns. $T_2$ viene
variato insieme a $T_1$ mantenendo $T_2=0.8T_1$.

\begin{table}[H]
\centering
\caption[Sweep di $T_1/T_2$]{TOP-4 ottiene $\bar M=1.00$ e successo
$100\%$ in tutti i punti.}
\label{tab:sweep_t1}
\begin{tabular}{ccccc}
\toprule
$T_1/T_2$ ($\mu$s) & $\bar M_1$ & $\bar M_{\mathrm{TOP4}}$ & $\rho$ & $p$ \\
\midrule
20/16 & 1.47 & 1.00 & 1.467 & $<0.001$ \\
50/40 & 2.20 & 1.00 & 2.200 & $<0.001$ \\
100/80 & 1.37 & 1.00 & 1.367 & $<0.001$ \\
200/160 & 1.60 & 1.00 & 1.600 & 0.001 \\
500/400 & 1.80 & 1.00 & 1.800 & $<0.001$ \\
\bottomrule
\end{tabular}
\end{table}

\begin{table}[H]
\centering
\caption[Sweep di $\lambda_{1q}$]{TOP-4 ottiene $\bar M=1.00$ e
successo $100\%$ in tutti i punti.}
\label{tab:sweep_eps1q}
\begin{tabular}{cccc}
\toprule
$\lambda_{1q}$ & $\bar M_1$ & $\rho$ & $p$ \\
\midrule
$10^{-4}$ & 1.60 & 1.600 & $<0.001$ \\
$5\times10^{-4}$ & 1.43 & 1.433 & $<0.001$ \\
$10^{-3}$ & 1.37 & 1.367 & $<0.001$ \\
$5\times10^{-3}$ & 1.67 & 1.667 & $<0.001$ \\
$10^{-2}$ & 1.57 & 1.567 & $<0.001$ \\
$2\times10^{-2}$ & 1.47 & 1.467 & $<0.001$ \\
\bottomrule
\end{tabular}
\end{table}

\begin{table}[H]
\centering
\caption[Sweep del readout simmetrico]{TOP-4 ottiene $\bar M=1.00$ e
successo $100\%$ in tutti i punti.}
\label{tab:sweep_pro}
\begin{tabular}{cccc}
\toprule
$p_{\mathrm{ro}}$ & $\bar M_1$ & $\rho$ & $p$ \\
\midrule
$0\%$ & 1.33 & 1.333 & 0.004 \\
$1\%$ & 1.50 & 1.500 & $<0.001$ \\
$2\%$ & 1.37 & 1.367 & $<0.001$ \\
$5\%$ & 1.47 & 1.467 & $<0.001$ \\
$10\%$ & 1.30 & 1.300 & 0.002 \\
$20\%$ & 1.47 & 1.467 & $<0.001$ \\
\bottomrule
\end{tabular}
\end{table}

Le oscillazioni non monotone di M1 sconsigliano inferenze causali sui
singoli parametri con questa numerosità. Il risultato supportato è più
limitato: TOP-4 non varia nei punti osservati.

% -----------------------------------------------
\subsection{Livello di Ottimizzazione}
\label{subsec:sensibilita_metodo}

\begin{table}[H]
\centering
\caption[Sweep di \texttt{optimization\_level}]{Conteggi della
compilazione effettiva e metriche su 30 repliche.}
\label{tab:sweep_optlevel}
\label{tab:sensitivity}
\begin{tabular}{cccccc}
\toprule
livello & CX & profondità & $\bar M_1$ & $\bar M_{\mathrm{TOP4}}$ & $\rho$ \\
\midrule
0 & 242 & 427 & 1.47 & 1.00 & 1.467 \\
1 & 236 & 421 & 1.50 & 1.00 & 1.500 \\
2 & 224 & 412 & 1.37 & 1.00 & 1.367 \\
3 & 224 & 412 & 1.37 & 1.00 & 1.367 \\
\bottomrule
\end{tabular}
\end{table}

% -----------------------------------------------
\subsection{Diagnostiche della struttura QPE}

Le diagnostiche decisive sono mantenute nel corpo: la
Figura~\ref{fig:istogrammi_qpe} confronta le repliche estreme,
l'Equazione~\ref{eq:miscela_coerente} definisce l'indice descrittivo e la
Tabella~\ref{tab:frazione_coerente} con la
Figura~\ref{fig:frazione_coerente} mostrano lo scarto dal proxy di nessun
evento. Gli sweep precedenti ne conservano qui il contesto parametrico.

% -----------------------------------------------
\subsection{Confronto con Zero-Noise Extrapolation}
\label{sec:app_zne}

La ZNE digitale scala $\lambda_{1q}$ e $\lambda_{2q}$, mantenendo fissi
coerenza e readout. ZNE-2 usa due livelli e ZNE-3 tre livelli condivisi;
non è gate folding su hardware.

\begin{table}[H]
\centering
\caption[Confronto ZNE appaiato]{I $p$ sono corretti con Holm nella
famiglia M1 contro ZNE-2, ZNE-3 e TOP-4.}
\label{tab:confronto_zne}
\begin{tabular}{lccccc}
\toprule
strategia & successo & $\bar M$ & shot/iter. & shot medi totali & $p_{\mathrm{Holm}}$ \\
\midrule
M1 & $100\%$ & 1.37 & 1\,024 & 1\,399 & --- \\
ZNE-2 & $100\%$ & 1.50 & 2\,048 & 3\,072 & 1.000 \\
ZNE-3 & $100\%$ & 1.57 & 3\,072 & 4\,813 & 1.000 \\
TOP-4 & $100\%$ & 1.00 & 1\,024 & 1\,024 & 0.0024 \\
\bottomrule
\end{tabular}
\end{table}

\paragraph{Validazione fuori dalla campagna rumorosa.}
I moltiplicatori Beauregard per $N=21/35$ sono validati idealmente e non
entrano negli sweep. Il contratto completo e il valore coerente con la
parametrizzazione Aer del canale a due qubit sono riportati nella
Sezione~\ref{app:fase1_contratto}; nessuno dei proxy è interpretato come
probabilità di successo.

\end{onehalfspace}
